% Part: turing-machines
% Chapter: machines-computations
% Section: halting-states

\documentclass[../../../include/open-logic-section]{subfiles}

\begin{document}

\olfileid{tur}{mac}{hal}
\olsection{حالات التوقف}

\begin{explain}
مع أننا عرّفنا آلاتنا بحيث لا تتوقف إلا عند عدم وجود تعليمة تنفذها،
فإن التمثيلات الشائعة لآلات تورنغ تشتمل على \emph{حالة توقف} مخصوصة
$h$، بحيث $h \in Q$.

والفكرة الكامنة وراء حالة التوقف بسيطة: حين تنتهي الآلة من عملها
(فتكون مستعدة لقبول الدخل، أو تكون قد فرغت من كتابة الخرج)، تنتقل إلى
حالة~$h$ وتتوقف فيها. ولبعض الآلات حالتا توقف، تقبل إحداهما الدخل
وترفضه الأخرى.
\end{explain}

\begin{ex}\emph{حالات التوقف}.
لإيضاح هذا المفهوم، لنبدأ بتعديل آلة الزوجية. فبدلًا من أن تتوقف
الآلة في الحالة~$q_0$ إذا كان الدخل زوجيًا، يمكننا إضافة تعليمة ترسل
الآلة إلى حالة توقف.
\[
\begin{tikzpicture}[->,>=stealth',shorten >=1pt,auto,node distance=2.8cm,
                    semithick]
  \tikzstyle{every state}=[fill=none,draw=black,text=black]

  \node[initial, state]         (A)                     {$q_0$};
  \node[state]         (B) [right of=A] {$q_1$};
  \node[state]         (C) [below of=A] {$h$};

  \path (A) edge [bend left] node {\TMtrans{\TMstroke}{\TMstroke}{R}} (B)
        edge node {\TMtrans{\TMblank}{\TMblank}{N}} (C)
        (B) edge [loop above] node {\TMtrans{\TMblank}{\TMblank}{R}} (B)
            edge [bend left] node {\TMtrans{\TMstroke}{\TMstroke}{R}} (A);
\end{tikzpicture}
\]

ولنوسع المثال أكثر. حين تقرر الآلة أن الدخل فردي، فإنها لا تتوقف قط.
ويمكننا تعديلها لتشتمل على حالة \emph{رفض}، وذلك باستبدال تعليمة
الدوران بتعليمة للانتقال إلى حالة رفض~$r$.
\[
\begin{tikzpicture}[->,>=stealth',shorten >=1pt,auto,node distance=2.8cm,
                    semithick]
  \tikzstyle{every state}=[fill=none,draw=black,text=black]

  \node[initial,state]         (A)                     {$q_0$};
  \node[state]         (B) [right of=A] {$q_1$};
  \node[state]         (C) [below of=A] {$h$};
  \node[state]         (D) [below of=B] {$r$};

  \path (A) edge [bend left] node {\TMtrans{\TMstroke}{\TMstroke}{R}} (B)
        edge node {\TMtrans{\TMblank}{\TMblank}{N}} (C)
        (B) edge node {\TMtrans{\TMblank}{\TMblank}{N}} (D)
            edge [bend left] node {\TMtrans{\TMstroke}{\TMstroke}{R}} (A);
\end{tikzpicture}
\]
\end{ex}

\begin{explain}
قد تفيد إضافة حالة توقف مخصوصة في حالات كهذه، إذ تبين صراحة متى تقبل
الآلة بعض المدخلات أو ترفضها. لكن من المهم ملاحظة أن إضافة حالة توقف
مخصوصة لا تمنح الآلة أي قدرة حسابية جديدة. وبالمثل، لمفهوم التوقف
الأقل صورية مزاياه أيضًا. فتعريف التوقف المستعمل حتى الآن في هذا
الفصل يجعل برهان \emph{مسألة التوقف} حدسيًا ويسير العرض. ولهذا سنواصل
استعمال تعريفنا الأصلي.
\end{explain}

\end{document}
